\subsection{Next-Phase Hazard Prediction Task (RQ2-i)}
\label{sec:prediction-task}

Part (i) of RQ2 asks how accurately next-phase team elimination can be
predicted from the features of Subsection~\ref{sec:methods-features} alone.
We call \(k\) the \textit{cut phase}: the model observes a match through the
end of phase \(k\) and predicts whether team \(i\) is eliminated during phase
\(k{+}1\). All inputs---features, edges, and masks---are constructed from
telemetry observable up to \(b_k\); events and outcomes after \(b_k\) enter
only the label definition.

Let \(b_k\) be the end time of phase \(k\), and \(T_i\) the elimination time
of team \(i\)---the moment its last surviving member is eliminated. The label
and hazard are
\[
Y_{i,k}=\Ind\!\left\{b_k < T_i \leq b_{k+1}\right\}, \qquad
h_{i,k}=P(b_k < T_i \leq b_{k+1}\mid T_i>b_k, X_{i,\leq k}).
\]
Teams eliminated before \(b_k\) leave the risk set and are masked. Surviving
teams contribute \(Y_{i,k}=0\) for every observed interval; once the match
ends or the eight-phase horizon is reached, their later intervals are
unobserved and treated as right-censored rather than labeled. Modeling uses
phases 1--8 and excludes the landing interval, so cut phases run from 1 to 7.
Training minimizes the masked discrete-time survival likelihood
\[
\mathcal{L}_{\mathrm{haz}}
= -\sum_{(m,i,k)\in \mathcal{R}}
\Big[ Y_{i,k}\log \hat h_{i,k} + (1-Y_{i,k})\log\big(1-\hat h_{i,k}\big) \Big],
\]
where the risk set \(\mathcal{R}\) contains exactly the team-phase cells with
\(T_i>b_k\) and an observed next phase; censored cells are excluded by the
mask. Splits are made by match, never by row. We answer part (i) in two ways:
by the absolute performance of feature-only models, and by a representation
ablation that degrades the feature set toward prior-style strategy scores
(Section~\ref{sec:experiments}).
